\documentclass[11pt]{article}

\usepackage[T1]{fontenc}
\usepackage{lmodern}
\usepackage[margin=1in]{geometry}
\usepackage{amsmath,amssymb}
\usepackage{booktabs}
\usepackage{microtype}
\usepackage{xcolor}
\pagecolor{white}
\usepackage{enumitem}
\definecolor{accent}{HTML}{142a5e}

\title{\textcolor{accent}{How Solid Is the Constraint-First Theorem?}\\[2mm]
\large every way we checked it, in plain English}
\author{Rob Mitchell}
\date{2026-10-23}

\begin{document}
\maketitle

\section*{The claim, in one sentence}

\textbf{A theory of everything cannot predict the numbers.} The only exact,
coordinate-independent structure in the string landscape is the D3 tadpole --- a single
``charge count'' constraint ($N = F^\top J H$). Everything else --- the couplings, the
masses, the cosmological constant --- is free input, not prediction. A theory of
everything predicts \emph{relations} (what is forbidden) and \emph{labels} (which
geometry), never \emph{values}.

That is a strong claim. This document lists, in plain English, every reason we believe it
--- and every way you could still prove us wrong.

\section*{Five independent ways we checked it}

\subsection*{1. The exact math proof (a computer did the algebra, to zero tolerance)}

We proved the core statement \emph{symbolically} --- not with approximate arithmetic, but
with exact algebra on fractions, so there is \emph{no rounding} anywhere. The result: the
tadpole is the \emph{unique} invariant, degree by degree, all five degrees checked. This
is a known theorem in mathematics (Weyl's First Fundamental Theorem, 1939); we re-derived
it and verified it by machine. One command reproduces the whole proof.

\subsection*{2. The machine-checked proof (a second computer checked every step)}

We wrote the proof again in \textbf{Lean 4}, a proof assistant whose kernel verifies
\emph{every single step} of the argument. There is no ``trust me'' --- the computer
confirms each step is valid, down to the metal. The result: degrees 1, 2, 3, and 4 of the
proof all build green (\emph{no \texttt{sorry}, no \texttt{axiom}, no gaps}). Two of the
hardest pieces (no cubic invariant, only $N^2$ at quartic) were checked this way.

\subsection*{3. The ``did we fool ourselves'' check (the definitions are real)}

A formal proof can lie if its definitions are secretly empty or over-constrained (a subtle
mistake can make every statement ``vacuously true''). We machine-checked that this did
\emph{not} happen: the symmetry group is non-empty and non-trivial (two distinct explicit
examples), the tadpole is \emph{not} identically zero (a concrete example gives $1$), and
the uniqueness theorem genuinely rules things out (the identity matrix, for one, is
\emph{not} an invariant). Five such sanity lemmas, all kernel-checked.

\subsection*{4. The measurements (real data, not just math)}

The math says the couplings should show no ``fingerprint'' of the geometry. We measured
that on real supersymmetric flux vacua: the mass spectrum's shape is statistically uniform
in flux (no fingerprint, as predicted), and the duality laws hold \emph{exactly} at
$h^{1,2}=2$ and $h^{1,2}=5$. The measurement agreed with the theorem.

\subsection*{5. The self-attack (independent reviewers tried to break it)}

We ran an \emph{adversarial} review: independent agents searched for a flaw. They found
one --- a counting error in a peripheral claim (``GUT collapses 3'' was wrong; it is 2).
We corrected it. The \emph{core} theorem survived every attack.

\section*{How you could still prove us wrong (the four falsifiers)}

A claim is only scientific if you can name what would break it. Here are the four, with an
honest sense of how likely each is.

\begin{center}
\begin{tabular}{lp{6.5cm}l}
\toprule
\# & falsifier & likelihood \\
\midrule
1 & a coordinate-independent quantity \emph{not} a function of the tadpole & $\approx$\,0\% (it is a proven theorem) \\
2 & an exact, derivable \emph{and} numerically-correct relation among the 19 (beyond gauge unification) & $\approx$\,30\% (the open frontier) \\
3 & a duality law that scatters at higher $h^{1,2}$ & $\approx$\,10\% (untested beyond $h^{1,2}=5$) \\
4 & a measured non-uniform distribution near $W_0=0$ & $\approx$\,5\% (measured flat) \\
\bottomrule
\end{tabular}
\end{center}

The key point: the most likely falsifier (\#2) attacks the \emph{conjectured extension} to
the Standard Model, \emph{not} the proven core. Even if it fires, the flux theorem stands.

\section*{The one mistake we caught (and fixed)}

The adversarial review caught a genuine error: we had said ``a grand-unified theory
collapses exactly 3 of the 19 parameters.'' It is actually 5 for the simplest such theory
(and up to 12 for others), and those extra relations are \emph{numerically wrong} against
data --- so only gauge unification (2) survives. We corrected every copy. The lesson: the
verification process worked exactly as intended --- it found a real mistake and the core
survived.

\section*{What we are still honest about (two gaps)}

\begin{enumerate}[leftmargin=1.4em]
\item \textbf{The full proof is machine-checked at degrees 1--4, not yet at every degree.}
  The all-degrees statement is a cited theorem (Weyl 1939) plus a degree-by-degree
  symbolic check (5 degrees); a full machine-checked proof at \emph{general} degree is a
  research project, not done.
\item \textbf{The Standard-Model extension is a conjecture, not a theorem.} We proved the
  19 parameters are free \emph{within} the Standard Model; that \emph{no} ultraviolet
  completion derives more than 2 of them correctly is the open conjecture (falsifier \#2).
\end{enumerate}

\section*{The bottom line}

The core of the theorem --- \emph{the only exact structure is the constraint} --- is as
solid as a mathematical claim can be: proven symbolically, machine-checked step by step,
checked for hidden emptiness, measured on real data, and stress-tested by an adversarial
review that caught its one real error. The remaining risk is concentrated in one place we
have always labeled open: the Standard-Model extension. Everything else is theorem, not
opinion.

\end{document}
